% Junior Level 3, Week 10 -- answer key
\answersweek{10}{Graduation Showcase}{Answers for workbook chapter 10 and play-at-home card 10.}

\section{Lesson questions}
\begin{answer}[Exit question]
Stop (hands off the keyboard), do not type it, and get a grown-up. The grown-up decides whether the site is OK; if it is, the child uses a nickname and the grown-up sets up any account. The same three steps work for a chatbot that asks where you live.
\end{answer}

\begin{answer}[Reflection]
Any drawing of a next idea and any honest smiley is right. Write the ideas down: games, stories and turtle pictures all lead into main-track Level 1.
\end{answer}

\begin{answer}[Show-me checklist]
The Week 10 show-me checklist is in the Assessments booklet. It is observed during the bug check, the typo rescue, the showcase and the gallery walk, and is never graded. \enquote{Getting there} is a normal result at this age.
\end{answer}

\section{Workbook exercises}
\begin{answer}[Exercise 10.1: Last-Minute Bug Check]
A finished project gets \enquote{yes} in all four rows. Common \enquote{not yet} answers and fixes: it starts differently the second time (Scratch: put \enquote{set \ldots{} to 0}, \enquote{erase all} or \enquote{show} under the green flag); a must-have is missing (add it, or point at it if it is there); hard to see (thicker pen, darker color). Fixing one thing is enough today.

\textbf{Typo rescue.} \texttt{t.forwad(100)}: the turtle window opens and the turtle stops at the line with the typo; the red message ends with \texttt{AttributeError} and says the turtle has no \texttt{forwad}, and may add \texttt{Did you mean: 'forward'?}. Fix: \texttt{forward}. A missing quote mark in \texttt{t.color("purple)}: nothing is drawn, and the message says \texttt{SyntaxError: unterminated string literal} with the line number. Fix: add the closing quote mark. A missing colon after \texttt{for i in range(4)}: nothing is drawn, and the message says \texttt{SyntaxError: expected ':'}. Fix: add the colon. \textbf{You should see after the fix:} an orange house with a purple roof, then the child's own color and number change. Look for: the child reads the line number and fixes only that line.
\end{answer}

\begin{answer}[Exercise 10.2: My Showcase]
Any true sentences about the child's own project are right. Example: \enquote{My plan card said a house, a roof and a sun. This part of my code draws the roof: it goes three times, because a triangle has three sides. I changed the house to blue. My idea came from the starter.} For Omar's loop: it runs three times; each time the turtle goes forward 100 and turns 120 (the turn-card number for a triangle). Look for: the child points at the right code and names one thing it does; a credit for any borrowed idea.
\end{answer}

\begin{answer}[Exercise 10.3: Sticker Feedback]
Any sticker placed with a kind spoken reason is right (\enquote{I liked the pattern, because it looks like a flower}). Look for: the visitor asks before trying, uses kind words only and does not change the project.
\end{answer}

\begin{answer}[A website asks for my name?]
The three steps: stop, do not type it, get a grown-up.
\end{answer}

\section{Play-at-home card}
\begin{answer}[Stop, do not type, get a grown-up]
A favorite color is fine to share. Full name and school: stop, do not type it, get a grown-up. A nickname for a score: ask a grown-up first; if they say yes, a nickname is fine. When the child is the website, praise tricky questions that the grown-up answers with the three steps.
\end{answer}

\begin{answer}[What next?]
Look for one part of the code explained in the child's own words. Any drawing of a next idea is right.
\end{answer}
